\section{Pre-registered LoRA configuration: falsification audit}
\label{supp:lora-audit}

\textit{Purpose of this supplement.}\;
This supplement documents the full audit chain for the LoRA
falsification arm, including the three pre-committed amendment
configurations, their collapse behaviour, and the three-criterion
failure rule that led to the methodological-null declaration.

This section documents the audit chain by which the parameter-efficient
fine-tuning arm was tested under three pre-committed amendments and
declared empirically undefined within the pre-registered compute
envelope. The point is methodological: a comparator that collapses to
the trivial all-negative solution is not a fair test of full
fine-tuning versus PEFT, regardless of which side ``wins'' in absolute
terms.

\subsection{Locked specification}

The pre-registered PEFT cell used LoRA \citep{hu2022lora} adapters
of rank 16 ($\alpha = 32$, dropout 0.05, no bias) on attention query
and value projections only --- with FFN, key projection, attention
output, embeddings, and pooler frozen. The classifier head was fully
trained (the \code{modules\_to\_save} parameter), giving approximately
0.541\% trainable parameters. Learning rate, budget, and all other
hyperparameters matched full fine-tuning ($2\times 10^{-5}$ for
2{,}048 updates).

\subsection{Three falsification attempts}

\paragraph{Attempt 1: locked specification.}
A 35-run pilot (PubMedBERT-base $\times$ T1B and T1F, 20 + 15 seeds)
collapsed identically: development macro-F1 stayed at the
all-\code{\_\_NEGATIVE\_\_} floor of $0.12624584\ldots$ across all 32
evaluation checkpoints per seed and across all 35 seeds. Test-set
per-head F1 was zero on every non-negative label across all 26
completed runs.

The internal evidence ruled out a wiring bug. LoRA adapter parameters
moved (max delta $8\times 10^{-3}$); the classifier head moved (max
delta $1.4\times 10^{-2}$); frozen parameters were strictly unchanged
(zero delta); training loss decreased monotonically from 2.20 to
0.20. The optimiser was working; the decision boundary simply never
left the trivial basin.

\paragraph{Attempt 2: 15$\times$ learning rate (retracted).}
Hypothesis: the locked LR was too low for a smaller trainable
footprint. A single-seed smoke at LR $3\times 10^{-4}$ (the
LoRA-paper default) produced bit-identical collapse. LoRA-B norm
grew 4--5$\times$ and the classifier weight norm grew 14\%,
confirming the optimiser had more signal --- the decision boundary
still did not move. The LR-cause hypothesis was falsified and this
amendment was retracted.

\paragraph{Attempt 3: 2$\times$ budget (falsified).}
Hypothesis: 2{,}048 updates were insufficient. Budget was doubled to
4{,}096 with LR restored to $2\times 10^{-5}$. Collapse persisted
across 64 evaluation checkpoints; training loss decreased 60\%; the
boundary still did not move.

\begin{table}[H]
\centering
\caption{Pre-committed acceptance criteria for the budget probe.}
\label{tab:s9-acceptance}
\small
\begin{tabular}{@{}lll@{}}
\toprule
\textbf{Criterion} & \textbf{Outcome} & \textbf{Pass?} \\
\midrule
Dev macro-F1 $> 0.20$ at any step $\le 1024$              & maximum 0.127 & no \\
Dev macro-F1 $> 0.30$ at step 4{,}096                        & 0.127         & no \\
At least one checkpoint with non-zero ex-NEG dev F1       & all zero      & no \\
\bottomrule
\end{tabular}
\end{table}

All three criteria failed. The pre-committed decision tree drops
the PEFT arm and reports the FT-vs-LoRA comparison as a
methodological null.

\subsection{Why this is not ``FT $>$ LoRA''}

Matched full fine-tuning escapes the all-negative basin within
approximately 512 steps and converges to dev macro-F1 of approximately
0.51. A naive read would conclude FT outperforms LoRA. We do not
draw this conclusion: the LoRA comparator is collapsed, not training.
Comparing a trained model to a collapsed comparator measures whether
the comparator collapses, not whether LoRA can compete with FT under
a richer configuration.

\subsection{What the audit rules out and rules in}

The collapse is a property of one cell's pre-registered
specification on this data: low trainable fraction (0.541\%),
attention-Q+V-only target modules, an 8-class head with $\sim 80\%$
\code{\_\_NEGATIVE\_\_} class imbalance, $\sim 1{,}200$ development
examples, and 2{,}048--4{,}096 optimiser steps. It does not falsify
LoRA in general or any LoRA result on different data, model size,
target-module set, or budget. A richer LoRA sweep (varying rank,
target modules beyond Q and V, custom warmup, class re-weighting)
would be needed to assess the LoRA-vs-FT comparison properly in
this small-data, high-imbalance regime; that study is reserved for
follow-up work.

\subsection{Disposition of degenerate runs}

\begin{table}[H]
\centering
\caption{Audit-archived degenerate runs. None enter any analysis,
figure, or table reported in the paper.}
\label{tab:s9-disposition}
\small
\begin{tabular}{@{}lr@{}}
\toprule
\textbf{Source} & \textbf{Runs} \\
\midrule
Attempt 1 (LR $2\!\times\!10^{-5}$, 2{,}048 steps)  & 35 \\
Attempt 2 (LR $3\!\times\!10^{-4}$, 2{,}048 steps)  & 1  \\
Attempt 3 (LR $2\!\times\!10^{-5}$, 4{,}096 steps)  & 1  \\
\midrule
\textbf{Total} & 37 \\
\bottomrule
\end{tabular}
\end{table}

\subsection{Realised configuration wave}

The pre-registered wave had been 3 encoders $\times$ 2 update regimes
$\times$ 3 schedules $\times$ 20 seeds = 360 main runs. Dropping the
PEFT regime leaves 3 $\times$ 1 $\times$ 3 $\times$ 20 = 180 main
runs plus 10 RoBERTa-base reference seeds. This is the realised
sample for every hypothesis test, variance decomposition, slope fit,
and rank-inversion quantification reported in the paper.